Table~\ref{tab:budget} reruns the two-district states New Hampshire, Maine and Rhode Island for budgets $s=1,\dots,4$ (the main regime is $s=1$). Brackets are closed under Proposition~\ref{prop:mono}(3): a plan that is feasible for $s$ is feasible for every larger budget, so the lower end at $s$ is the largest lower end over budgets $\le s$, and the upper end is the smallest over budgets $\ge s$ and the geography-free bound.

\begin{table}[t]
\centering\small
\caption{Best-district ($q=1$) and second-best ($q=2$) share, in percent of the two-party vote, as the county-split budget $s$ grows ($\pm1\%$, 2020 presidential vote); geography-free hull bound and unconstrained ReCom incumbent for comparison. ``$<$50'': no plan of the regime has $q$ districts at 50\% or more.}\label{tab:budget}
\begin{tabular}{llcccccrr}
\toprule
State & Party & $q$ & $s=1$ & $s=2$ & $s=3$ & $s=4$ & geo-free & ReCom (no budget) \\
\midrule
New Hampshire & D & 1 & 57.02--57.5 & 57.02--60.5 & 57.02--61.5 & 57.02--61.5 & 62.4 & 57.0 \\
New Hampshire & D & 2 & 53.75--53.8 & 53.75--53.8 & 53.75--53.8 & 53.75--53.8 & 53.8 & 53.7 \\
New Hampshire & R & 1 & $<$50 & 50.01--52.5 & 50.01--54.6 & 50.01--54.6 & 54.6 & -- \\
New Hampshire & R & 2 & $<$50 & $<$50 & $<$50 & $<$50 & $<$50 & -- \\
Maine & D & 1 & 62.67--63.0 & 63.87--64.0 & 64.07--64.5 & 64.51--65.0 & 66.7 & 63.3 \\
Maine & D & 2 & 54.62--54.7 & 54.62--54.7 & 54.62--54.7 & 54.62--54.7 & 54.7 & 54.6 \\
Maine & R & 1 & 54.52--55.0 & 55.54--56.0 & 55.82--56.0 & 56.02--56.5 & 57.9 & 54.8 \\
Maine & R & 2 & $<$50 & $<$50 & $<$50 & $<$50 & $<$50 & -- \\
Rhode Island & D & 1 & 70.04--70.5 & 70.04--71.5 & 70.37--72.4 & 70.37--72.4 & 72.4 & 69.6 \\
Rhode Island & D & 2 & 60.60--60.6 & 60.60--60.6 & 60.60--60.6 & 60.60--60.6 & 60.6 & 60.6 \\
Rhode Island & R & 1 & $<$50 & $<$50 & $<$50 & $<$50 & $<$50 & -- \\
Rhode Island & R & 2 & $<$50 & $<$50 & $<$50 & $<$50 & $<$50 & -- \\
Idaho & D & 1 & $<$50 & $<$50 & $<$50 & $<$50 & $<$50 & -- \\
Idaho & D & 2 & $<$50 & $<$50 & $<$50 & $<$50 & $<$50 & -- \\
Idaho & R & 1 & 75.50--76.0 & 75.88--77.0 & 75.88--78.0 & 75.88--78.0 & 79.6 & 74.5 \\
Idaho & R & 2 & 65.88--65.9 & 65.88--65.9 & 65.88--65.9 & 65.88--65.9 & 65.9 & 65.9 \\
Montana & D & 1 & 50.00--52.5 & 50.00--57.3 & 50.00--57.3 & 50.00--57.3 & 57.3 & -- \\
Montana & D & 2 & $<$50 & $<$50 & $<$50 & $<$50 & $<$50 & -- \\
Montana & R & 1 & 67.53--69.5 & 67.53--71.0 & 68.38--71.0 & 68.38--72.9 & 72.9 & 66.6 \\
Montana & R & 2 & 58.40--58.4 & 58.40--58.4 & 58.40--58.4 & 58.40--58.4 & 58.4 & 58.4 \\
West Virginia & D & 1 & $<$50 & $<$50 & $<$50 & $<$50 & $<$50 & -- \\
West Virginia & D & 2 & $<$50 & $<$50 & $<$50 & $<$50 & $<$50 & -- \\
West Virginia & R & 1 & 76.00--77.0 & 76.00--77.5 & 76.00--79.0 & 76.00--79.5 & 80.4 & 74.0 \\
West Virginia & R & 2 & 69.78--69.8 & 69.78--69.8 & 69.79--69.8 & 69.79--69.8 & 69.8 & 69.8 \\
\bottomrule
\end{tabular}

\end{table}

\paragraph{Capacity grows with every split, provably.}
For Maine the certified best-district share of Democrats is 62.67--63.0\% at $s=1$, 63.87--64.0\% at $s=2$, 64.07--64.5\% at $s=3$ and 64.51--65.0\% at $s=4$ (Republicans: 54.5--55.0, 55.5--56.0, 55.8--56.0, 56.0--56.5\%). Wherever the upper end of one budget lies below the lower end of the next (Maine $s=1\to2$ for both parties, Democrats $s=2\to3$ and $3\to4$, Republicans $3\to4$), the extra split is \emph{certified} to increase capacity, by at least 0.8 point for the first split of Maine Democrats. The geography-free bound (66.7\% and 57.9\%) is far above even $s=4$: contiguity and the remaining budget, not the vote totals alone, bind in these states.

\paragraph{A seat-level effect.}
New Hampshire Republicans cannot make a single 50\% district with one split (certified), but can with two: a verified plan has a 50.01\% district. The frontier therefore moves by a whole seat at the first extra split.

\paragraph{The all-safe end does not depend on the budget.}
$\sigma^*(2)$ is pinned to the statewide share for every budget tried (New Hampshire 53.75--53.8\%, Maine 54.62--54.7\%, Rhode Island 60.60\%), as it must be when the largest possible value equals the statewide margin (Section~\ref{sec:main}).

\paragraph{Tractability.}
Exactness is lost when the budget grows in New Hampshire (upper end for Democrats 57.5\% at $s=1$, 60.5\% at $s=2$, 61.5\% at $s\ge3$) and Rhode Island, but not in Maine, where every bracket is at most half a point wide for $s\le4$. This is the tractability frontier of the method: it depends on the budget and on the structure of the counties (Maine has 16), and it is where a relaxed solution has room to cheat inside a split county.
